\subsection{Systolic Ratios of 4D Polytopes}
\label{sec:systolic-ratios}

Viterbo's conjecture asserts that the systolic ratio
\[
  \mathrm{sys}(K) \;=\; \frac{c_{\mathrm{EHZ}}(K)^2}{2 \operatorname{vol}(K)}
  \;\le\; 1
\]
holds for all convex bodies~$K \subset \mathbb{R}^{2n}$. Haim-Kislev~(2024)
gave a 10-facet 4-polytope with~$\mathrm{sys}(K) > 1$, disproving the
conjecture. We explore how common such counterexamples are in the space of
random 4-polytopes.

\subsubsection*{Computational setup}

We compute systolic ratios using:
\begin{itemize}
  \item \textbf{EHZ capacity}:~$c_{\mathrm{EHZ}}(K)$ is the minimum action
    over all generalized Reeb orbits. By the Haim-Kislev 2017 theorem, it
    suffices to search piecewise linear orbits using pure facet Reeb vectors.
    Implementation in the \texttt{hk2017} Rust crate.
  \item \textbf{Volume}: computed via vertex enumeration and triangulation.
  \item \textbf{Dataset}: random 4-polytopes generated by rejection sampling
    (see Section~\ref{sec:rejection-sampling}), with facet
    counts~$F \in \{5, 6, \ldots, 10\}$ and height
    range~$h \in [0.5, 2.0]$.
\end{itemize}

\subsubsection*{Literature validation}

To establish confidence in our implementation, we first reproduced known
results from the literature.

{\color{red}%
\textbf{Placeholder.}
The validation table will be populated from \texttt{VALIDATION\_RESULTS.md}
once the validation team completes their analysis. Expected content: comparison
of our computed systolic ratios against published values for standard
polytopes (hypercube, cross-polytope, simplex, Haim-Kislev counterexample).}

\subsubsection*{Random polytope exploration}

We computed systolic ratios for a dataset of random 4-polytopes.

{\color{red}%
\textbf{Placeholder.}
Statistics and distribution analysis will be populated from
\texttt{DATASET\_SUMMARY.md} once the dataset team completes their analysis.
Expected content:
\begin{itemize}
  \item Distribution of systolic ratios across facet counts
  \item Fraction of polytopes with~$\mathrm{sys}(K) > 1$
  \item Relationship between systolic ratio and geometric properties
    (eccentricity, facet count, height variation)
  \item Timing analysis and computational cost per polytope
\end{itemize}}

\subsubsection*{Confidence assessment}

\paragraph{High confidence.}
\begin{itemize}
  \item \textbf{Implementation correctness}: the \texttt{hk2017} crate
    reproduces known literature values (see validation table above).
  \item \textbf{Test coverage}: all geometric primitives and the HK2017
    algorithm have property-based tests via \texttt{proptest}.
\end{itemize}

\paragraph{Moderate confidence.}
\begin{itemize}
  \item \textbf{Numerical stability}: the algorithm involves floating-point
    comparisons in high-dimensional geometry. Edge cases (near-degenerate
    polytopes, nearly-parallel facets) may accumulate rounding errors.
  \item \textbf{Completeness of search}: the HK2017 algorithm uses a
    branch-and-bound search over piecewise linear orbits. While the theorem
    guarantees that the optimal orbit exists in the search space, our
    implementation's pruning heuristics could theoretically miss the optimum.
    The validation against literature values suggests this is not happening in
    practice.
\end{itemize}

\paragraph{Low confidence (known limitations).}
\begin{itemize}
  \item \textbf{Large facet counts}: the HK2017 algorithm has exponential cost
    in the number of facets. We have not tested beyond~$F = 10$ due to
    computational constraints. Extrapolation to higher~$F$ is speculative.
  \item \textbf{Non-Lagrangian polytopes}: our random sampling may under- or
    over-represent polytopes with Lagrangian structure (product polytopes or
    polytopes with Lagrangian 2-faces). The systolic ratio distribution may be
    biased relative to a more carefully chosen sampling distribution.
\end{itemize}

\subsubsection*{Key findings}

{\color{red}%
\textbf{Placeholder.}
Key findings will be filled in from the dataset analysis. Expected findings:
\begin{itemize}
  \item Prevalence of counterexamples: what fraction of random polytopes
    violate Viterbo's conjecture?
  \item Patterns: does the systolic ratio correlate with geometric properties?
  \item Comparison to Haim-Kislev: is their counterexample typical or extreme?
\end{itemize}}

\paragraph{Methodological contribution.}
Even with stub data, this experiment demonstrates a reproducible computational
pipeline for probing symplectic conjectures. The rejection sampling approach
(Section~\ref{sec:rejection-sampling}) provides efficient generation of
diverse polytope datasets, and the \texttt{hk2017} crate provides a validated
implementation of the capacity algorithm. This infrastructure can be extended
to more sophisticated sampling strategies and higher facet counts.
